% eksperymentalne tłumaczenie na włoski

%

TODO: Punkt Nagela, pojawia się u Evesa \cite[s. 71]{eves1_1972} jako ćwiczenie na zastosowanie twierdzenia Cevy razem z punktem Gergonne'a. % lang-pl

\index{punkt!Nagela} % lang-pl
TODO: Il punto di Nagel compare in Eves \cite[p. 71]{eves1_1972} come esercizio sull'applicazione del teorema di Ceva insieme al punto di Gergonne. % lang-it

\index{czewiany Nagela} % lang-pl
\index{punto!di Nagel} % lang-it

\index{ceviane di Nagel} % lang-it

\begin{proposition}[punkt Nagela] % lang-pl
\begin{proposition}[punto di Nagel] % lang-it
    \label{punkt_nagela}
\end{proposition}

O punkcie Nagela pisze też Zetel \cite[s. 22, 25, 64]{zetel_2020} (punkt Nagela, punkt przecięcia środkowych -- środek ciężkości trójkąta, środek okręgu wpisanego oraz środek ciężkości obwodu trójkąta leżą na jednej prostej zwanej drugą prostą Eulera albo prostą Eulera-Nagela). % lang-pl

Del punto di Nagel scrive anche Zetel \cite[p. 22, 25, 64]{zetel_2020} (punto di Nagel, punto di concorrenza delle mediane: il baricentro del triangolo, l'incentro e il baricentro del perimetro del triangolo giacciono su una stessa retta, detta seconda retta di Eulero oppure retta di Eulero--Nagel). % lang-it

%